\documentclass[main]{subfiles}
\usepackage[]{workbook}
\numbering{ws2-7}

\begin{document}

\ifSubfilesClassLoaded{%
  \chaptermark{\chaptertwo}
}{}

\section{The Direct Comparison Test} \label{ws2-7}
\begin{framed}
\centerline{\textbf{Lesson Goals}}

You should be able to:
\begin{itemize}
\item Apply the Direct Comparison Test to determine if a given series converges or diverges
\item Decide the best strategy for comparing a given series to a known series, as well as determine when a comparison is inconclusive, and when a comparison cannot be made
\item Use the Direct Comparison Test to reason abstractly about the convergence of a series 
\end{itemize}
\end{framed}

 \begin{enumerate}

  
   \item
We saw before that the \textit{harmonic series} $\displaystyle \sum_{k=1}^{\infty} \frac1k = 1+\frac12 +\frac13 +\cdots$ diverges, by comparing it to the improper integral $\displaystyle \int_1^{\infty} \frac1x~dx$. However, the divergence of the harmonic series was known since the 1300's, centuries before the discovery of calculus.

Consider the series $\displaystyle \sum_{k=1}^{\infty} a_k = \underbrace{\frac{1}{2}}_{\textrm{1 term}} + \underbrace{\frac{1}{4} +
      \frac{1}{4}}_{\textrm{2 terms}} + \underbrace{\frac{1}{8} + \cdots +
      \frac{1}{8}}_{\textrm{$4$ terms}} + \underbrace{\frac{1}{16} +
      \cdots + \frac{1}{16}}_{\textrm{$8$ terms}} +
      \underbrace{\frac{1}{32} + \cdots + \frac{1}{32}}_{\textrm{$16$
      terms}} + \cdots$.

    \begin{enumerate}
      \item What does the Nth Term Test for Divergence say about the series $\displaystyle \sum_{k=1}^{\infty} a_k$?
      \begin{problem}[0.5in]
      
      \end{problem}

      \begin{solution}
      Nothing.  $\displaystyle \lim_{k \to \infty} a_k = 0$, so
      the Nth Term Test for Divergence gives no information about the convergence of this series.
      \end{solution}

      \item    
      \begin{problem}[1in]
      Does the series $\displaystyle \sum_{k=1}^{\infty} a_k$ converge or diverge? Why?
      \end{problem}
      
      
      \begin{solution}
      Notice that each set of bracketed terms sums to $\frac12$, which means that each pair of bracketed terms sums to 1. Since there are infinitely many sets of bracketed terms, we can find a partial sum that exceeds any number, by simply adding enough terms until we've added on another 1 from the bracketed terms. More explicitly, the partial sum with $2^n-1$ terms is equal to  $\frac{n}{2}$; by making $n$ as large as we want, we can make the partial sums as large as we want as well.  Thus, the partial sums of the series tends to $\infty$, so by the Definition of Convergence the series \fbox{diverges}.   
      \end{solution}
      
      \item
      \begin{problem}
      How do the individual terms of the harmonic series compare to the corresponding terms in the series $\displaystyle \sum_{k=1}^{\infty} a_k$? What does this tell us about the harmonic series?
      \end{problem}

      \begin{solution}
      Let's write out the terms of the harmonic series on top of the terms for the series $\displaystyle \sum_{k=1}^{\infty} a_k$:
      
      \begin{align*}
      \displaystyle \sum_{k=1}^{\infty} \frac1k &= 1+\underbrace{\frac12 +\frac13} +\underbrace{\frac14 +\frac15 +\frac16 +\frac{1}{7}}+\underbrace{\frac{1}{8}+\frac{1}{9}+\frac{1}{10}+\frac{1}{11}+\frac{1}{12}+\frac{1}{13}+\frac{1}{14}+\frac{1}{15}}+ \cdots\\
      \displaystyle \sum_{k=1}^{\infty} a_k &= \frac{1}{2} + \underbrace{\frac{1}{4} +
      \frac{1}{4}} + \underbrace{\frac{1}{8} + \frac{1}{8}+\frac{1}{8}+
      \frac{1}{8}} + \underbrace{\frac{1}{16} + \frac{1}{16}+\frac{1}{16} + \frac{1}{16} + \frac{1}{16} + \frac{1}{16} + \frac{1}{16} + \frac{1}{16}} + \cdots
      \end{align*}
      
      Notice that each term of the harmonic series is larger than its corresponding term in the second series. Also notice that all of the terms in both series are positive. Thus, we can say that:
      
      $$0<a_k< \frac{1}{k}$$
      for all $k\geq 1$. 
      
      Therefore, every partial sum of $\displaystyle \sum_{k=1}^{\infty} \frac1k$ is larger than the corresponding partial sum from $\displaystyle \sum_{k=1}^{\infty} a_k$. Since we know that the partial sums of the second series approach infinity, the partial sums from the harmonic series must approach infinity as well. So once again, we see that the harmonic series must diverge, \textit{by direct comparison to another series}.
      
      It is important
      to write the $0 < a_k$ part of the inequality because it says that
      the terms of both series are all positive, which is why the
      Direct Comparison Test can be used.
      \end{solution}

    \end{enumerate}

\pspace{\newpage}

\begin{framed}
\centerline{\textbf{The Direct Comparison Test}}
Let $\displaystyle \sum a_k$ and $\displaystyle \sum b_k$ be two series such that 
$$0\leq a_k \leq b_k$$ for all $k$, \textit{eventually}. Then:
\begin{itemize}
\item If $\displaystyle \sum b_k$ \uline{converges}, then $\displaystyle \sum a_k$ \uline{converges} as well. 
\item If $\displaystyle \sum a_k$ \uline{diverges}, then $\displaystyle \sum b_k$ \uline{diverges} as well. 
\end{itemize}

\textcolor{red}{\uline{Warning}: If $\displaystyle \sum a_k$ \textit{converges}, we cannot say anything about the convergence of $\displaystyle \sum b_k$! Similarly, if $\displaystyle \sum b_k$ \textit{diverges}, we cannot say anything about the convergence of $\displaystyle \sum a_k$!}
\end{framed}



 \item  Use the Direct Comparison Test to determine whether each of the following series converges or diverges.
 
    \probonly{}
    \begin{enumerate}
     \item
      \begin{problem}[1in]
      $\displaystyle \sum_{k=1}^\infty \frac{|\cos k|}{k^3}$.
      \end{problem}

      \begin{solution}
      We know that $0 \leq |\cos k| \leq 1$ for all $k$, so $\displaystyle
      0 \leq \frac{|\cos k|}{k^3} \leq \frac{1}{k^3}$.  Since
      $\displaystyle \sum_{k=1}^\infty \frac{1}{k^3}$ converges (it's a
      $p$-series with $p = 3>1$), the
      Direct Comparison Test shows that $\displaystyle \sum_{k=1}^\infty
      \frac{|\cos k|}{k^3}$ \fbox{converges}.
      \end{solution}
      
       \item 
    \begin{problem}[1in]
    $\displaystyle \sum_{k=1}^\infty \frac{-1}{2^k + k}$
    \end{problem}

    \begin{solution}
    First, notice that all of the terms of this series are negative, so we can't use the Direct Comparison Test on it immediately. However, we can simply factor the $-1$ out of all the terms, and work with the positive series $\displaystyle \sum_{k=1}^\infty \frac{1}{2^k + k}$, because multiplying a series by $-1$ does not affect its convergence/divergence.
    
    We know that $0 < \frac{1}{2^k + k} < \frac{1}{2^k}$ for all $k \geq
    1$.   We also know that $\displaystyle \sum_{k=1}^\infty
    \frac{1}{2^k}$ converges, since it is a geometric series with $|r|=\frac12<1$. (We initially showed that this series converges by showing that the limit of its partial sums is finite.) So, the Direct Comparison Test says that
    $\displaystyle \sum_{k=1}^\infty \frac{1}{2^k + k}$ converges, too.
    
    Therefore, the original series $\displaystyle \sum_{k=1}^\infty \frac{-1}{2^k + k}$ \fbox{converges}.
    \end{solution}



      \item 
      \begin{problem}[1in]
      $\displaystyle \sum_{k=100}^\infty \frac{5 + 3\sin k}{k}$.
      \end{problem}

      \begin{solution}
      Here (in \ideacolor) is how you might get some intuition about
      whether the series converges: \ideas{We are basically interested in
      how quickly the terms of the series go to 0.  In this case, the
      numerator $5 + 3 \sin k$ oscillates between $5 + 3 (-1) = 2$ and
      $5 + 3(1) = 8$, so it doesn't have much impact on how quickly the
      terms of the series go to 0.  Therefore, we can think of the terms
      of the series as being ``similar'' to $\frac{1}{k}$, so it seems
      like the series should diverge.}

      To justify our intuition mathematically, we use the Direct Comparison Test.
      Since we believe the series diverges, we need to compare with a
      series whose terms are smaller than the terms of our series.  To
      come up with such a series, we use the fact that $5 + 3 \sin k \geq
      5 + 3 (-1) = 2$, no matter what $k$ is.  So, here is a complete
      argument:

      For all $k>0$, $\displaystyle 0\leq  \frac{2}{k}\leq \frac{5 + 3 \sin k}{k}$.  We know that $\displaystyle \sum \frac{2}{k}$
      diverges\footnote{When we write $\sum a_k$ without any bounds of
      summation, we are always referring to an infinite series, but we
      are not precise about where exactly the series starts.  That is,
      this could mean $\displaystyle \sum_{k=10}^\infty a_k$ or
      $\displaystyle \sum_{k=427}^\infty a_k$ (or something else).
      Since we know that the beginning terms of a series don't affect
      convergence, it's not necessary to be particularly precise about
      where our series starts when we're just discussing whether it
      converges.}, since it's simply a constant multiple of the divergent harmonic series. So the Direct Comparison
      Test tells us that $\displaystyle \sum \frac{5 + 3 \sin k}{k}$ also
      \fbox{diverges}.
      \end{solution}
      
      \item 
      \begin{problem}[1in]
      $\displaystyle \sum_{k=37}^\infty \frac{5 + 3\sin k}{2^k}$.
      \end{problem}

      \begin{solution}
      Here, the same type of intuitive reasoning as in the previous part
      leads us to guess that the series converges.  So, we need to think
      of a larger series to compare with.  Here is a complete argument:

      For all $k>0$, $\displaystyle 0 < \frac{5 + 3 \sin k}{2^k} \leq
      \frac{5 + 3(1)}{2^k} = \frac{8}{2^k}$.  We know $\displaystyle \sum \frac{8}{2^k}$
      converges, since it's 8 times the series $\displaystyle \sum \frac{1}{2^k}$, which we know is a convergent geometric series.  So, the Direct Comparison Test says that
      $\displaystyle \sum \frac{5 + 3\sin k}{2^k}$ \fbox{converges}, too. 
      \end{solution}
      
            \end{enumerate}       
\pspace{\newpage}

\item Determine if the following series converge or diverge.
\begin{enumerate}
  \item 
    \begin{problem}[1in]
    $\displaystyle \sum_{k = 10^{10}}^\infty \frac{1}{100000000k}$
        \end{problem}

    \begin{solution}
    We know that the harmonic series $\displaystyle \sum_{k=1}^\infty
    \frac{1}{k}$ diverges.  Therefore, the series $\displaystyle
    \sum_{k=1}^\infty \frac{1}{100000000k}$ diverges, since it's just the
    harmonic series multiplied by a constant.  We know that the beginning
    terms of a series do not affect its convergence, so if we start with
    the $10^{10}$-th term, the series still \fbox{diverges}.
    \end{solution}

   
    \item
      \begin{problem}[1in]
      $\displaystyle \sum_{k=1}^\infty \frac{3 + \sin k}{0.9^k}$.
      \end{problem}

      \begin{solution}
      We know that $-1 \leq \sin k \leq 1$ for all $k$, so $2 \leq 3 +
      \sin k \leq 4$ for all $k$.  Therefore, $\displaystyle 0 \leq
      \frac{2}{0.9^k} \leq \frac{3 + \sin k}{0.9^k} \leq \frac{4}{0.9^k}$
      for all $k$.  We know that $\displaystyle \sum_{k=1}^\infty
      \frac{2}{0.9^k}$ and $\displaystyle \sum_{k=1}^\infty
      \frac{4}{0.9^k}$ both diverge, since they are both geometric series with
      common ratio $r = \frac{1}{0.9}$, and geometric series with $|r|
      \geq 1$ diverge.  The relevant inequality is the fact that the
      series in question is greater than something divergent, so here is a
      complete argument.

      We know $\displaystyle 0 \leq \frac{2}{0.9^k} \leq \frac{3 + \sin
      k}{0.9^k}$ for all $k$, and we've also explained why
      $\displaystyle \sum_{k=1}^\infty \frac{2}{0.9^k}$ diverges.
      Therefore, the Direct Comparison Test shows that $\displaystyle
      \sum_{k=1}^\infty \frac{3 + \sin k}{0.9^k}$ \fbox{diverges}.
      
      However, we could have also concluded divergence using the Nth Term Test for Divergence, as $$\lim_{k \to \infty} \frac{3 + \sin k}{0.9^k} = \lim_{k \to \infty} \frac{3 + \sin k}{\left(\frac{9}{10}\right)^k} = \lim_{k \to \infty} \left(\frac{10}{9}\right)^k (3 + \sin k) =\infty \not= 0$$
     
     This again demonstrates the usefulness of first checking if the terms of a series go to 0 before applying other strategies for determining convergence, as it is often easier to apply than the Definition of Convergence or the Comparison Tests.
      \end{solution}
\end{enumerate}






\item \label{spring2007-midterm2-3-modified}

  
  Suppose $\displaystyle \sum_{k=1}^{\infty} a_k$ converges, and $a_k > 0$
  for all $k$.  Which of the following statements must be true, which must be false, and for which
  is it impossible to decide?  Explain your reasoning clearly, and cite
  any convergence tests/strategies that you use.

  \begin{enumerate}
  \item
    \begin{problem}[\vfill]
      $\displaystyle\sum_{k=1000}^\infty \frac{a_k}{k}$ converges.        
    \end{problem}
    
    \begin{solution}
      \fbox{Must be true.}  We can use direct comparison.  For $k \geq
      1000$, $0 < \frac{a_k}{k} < a_k$.  Since $\sum a_k$ converges, the
      \textbf{Direct Comparison Test} says that $\sum \frac{a_k}{k}$
      converges, too.
    \end{solution}
    
    
    
    
     \item
    \begin{problem}[\vfill]
      $\displaystyle\sum_{k=1}^{\infty} (a_k)^2$ converges.         
    \end{problem}

    \begin{solution}

      \fbox{Must be true.}  Since $\sum a_k$ converges, the Nth Term Test
      tells us that $\displaystyle \lim_{k \to \infty} a_k = 0$.  Since
      $a_k$ is always positive, this tells us that, eventually (for $k$
      large enough) $a_k$ is always less than 1.  When $a_k < 1$, $0 <
      (a_k)^2 < a_k$, so we can use the \textbf{Direct Comparison Test}:
      $\sum a_k$ converges, so $\sum (a_k)^2$ does as well.
    \end{solution}





  \item \label{spring2007-midterm2-3-modified-1}
    \begin{problem}[\vfill]
      $\displaystyle\sum_{k=1}^{\infty} \frac{1}{a_k}$ converges.  
    \end{problem}

    \begin{solution} 
      If you're not sure whether such a statement is true, it can be
      helpful to use our key intuitive idea that a series converges if its
      terms go to 0 quickly enough.  For instance, here (in maroon) is
      some intuitive reasoning along those lines: \ideas{Since we are told
        that $\sum a_k$ converges, we know that $a_k \to 0$ quickly.  We
        are wondering whether $\sum \frac{1}{a_k}$ converges, so we would
        like to know whether $\frac{1}{a_k} \to 0$ quickly.  But since
        $a_k \to 0$, $\frac{1}{a_k}$ does not tend to 0 at all, so $\sum
        \frac{1}{a_k}$ must diverge.}  Once you have the intuitive idea,
      use that to figure out what definition or convergence test you can
      use to justify your reasoning mathematically.  In this case, the
      heart of our intuitive argument was that $\frac{1}{a_k}$ does not
      tend to 0 at all, which is basically the Nth Term Test.  So, here is
      a solution with complete mathematical justification:

      \fbox{Must be false}.  Since $\displaystyle \sum_{k=1}^\infty a_k$
      converges, we know that $a_k \to 0$ as $k \to \infty$ (this is a
      consequence of the Nth Term Test).  Therefore, as $k \to \infty$,
      $\frac{1}{a_k} \to +\infty$ (since $a_k >0$ for all $k$).  By the \textbf{Nth Term Test}, since $\displaystyle \lim_{k \to \infty}a_k \not=0$,
      $\displaystyle \sum_{k=1}^\infty \frac{1}{a_k}$ must diverge.
    \end{solution}
    
      
   \item
    \begin{problem}[\vfill]
      $\displaystyle\sum_{k=1}^{\infty} \ln a_k$ converges.
    \end{problem}

    \begin{solution}
      \fbox{Must be false.}  Since $\sum a_k$ converges, the Nth Term Test
      tells us that $\displaystyle \lim_{k \to \infty} a_k = 0$.
      Therefore, $\displaystyle \lim_{k \to \infty} \ln a_k = - \infty$.
      Since this limit is not equal to 0, the \textbf{Nth Term Test} tells
      us that $\sum \ln a_k$ must diverge.
    \end{solution}
  \end{enumerate}

  \begin{solutiononly} {\LARGE\Pointinghand} As mentioned in
    \ref{spring2007-midterm2-3-modified-1}, it's often helpful in problems
    like this to appeal to our key intuitive idea that a series converges
    if its terms go to 0 quickly enough.
    It's also a good idea to check your intuition using some series you
    know, like $p$-series or geometric series.

    Remember, though, that NONE of that constitutes mathematical
    justification; to justify your answer, you'll still need to use a
    convergence test or definition!
  \end{solutiononly}

  

  
  
 
      
      
       \probonly{\newpage}

    % \item  Decide whether each of the following series converges, and use the Direct
    % Comparison Test to justify your answer mathematically.

 
  
    % \begin{enumerate}
    %   \item 
    %   \begin{problem}
    %   $\displaystyle \sum_{k=100}^\infty \frac{5 + 3\sin k}{k}$.
    %   \end{problem}

    %   \begin{solution}
    %   Here (in \ideacolor) is how you might get some intuition about
    %   whether the series converges: \ideas{We are basically interested in
    %   how quickly the terms of the series go to 0.  In this case, the
    %   numerator $5 + 3 \sin k$ oscillates between $5 + 3 (-1) = 2$ and
    %   $5 + 3(1) = 8$, so it doesn't have much impact on how quickly the
    %   terms of the series go to 0.  Therefore, we can think of the terms
    %   of the series as being ``similar'' to $\frac{1}{k}$, so it seems
    %   like the series should diverge.}

    %   To justify our intuition mathematically, we use the Comparison Test.
    %   Since we believe the series diverges, we need to compare with a
    %   series whose terms are smaller than the terms of our series.  To
    %   come up with such a series, we use the fact that $5 + 3 \sin k \geq
    %   5 + 3 (-1) = 2$, no matter what $k$ is.  So, here is a complete
    %   argument:

    %   For all $k>0$, $\displaystyle \frac{5 + 3 \sin k}{k} \geq \frac{2}{k}
    %   > 0$.  We know that $\displaystyle \sum \frac{2}{k}$
    %   diverges\footnote{When we write $\sum a_k$ without any bounds of
    %   summation, we are always referring to an infinite series, but we
    %   are not precise about where exactly the series starts.  That is,
    %   this could mean $\displaystyle \sum_{k=10}^\infty a_k$ or
    %   $\displaystyle \sum_{k=427}^\infty a_k$ (or something else).
    %   Since we know that the beginning terms of a series don't affect
    %   convergence, it's not necessary to be particularly precise about
    %   where our series starts when we're just discussing whether it
    %   converges.} (it's twice the harmonic series), so the Comparison
    %   Test tells us that $\displaystyle \sum \frac{5 + 3 \sin k}{k}$ also
    %   diverges.
    %   \end{solution}
      
    %     \probonly{\vfill}
    
      
    %   \item 
    %   \begin{problem}
    %   $\displaystyle \sum_{k=37}^\infty \frac{5 + 3\sin k}{2^k}$.
    %   \end{problem}

    %   \begin{solution}
    %   Here, the same type of intuitive reasoning as in the previous part
    %   leads us to guess that the series converges.  So, we need to think
    %   of a larger series to compare with.  Here is a complete argument:

    %   For all $k>0$, $\displaystyle 0 < \frac{5 + 3 \sin k}{2^k} \leq
    %   \frac{5 + 3(1)}{2^k} = \frac{8}{2^k}$.  We know $\displaystyle \sum \frac{8}{2^k}$
    %   converges (it's 8 times the series $\sum \frac{1}{2^k}$, and we know
    %   $\sum \frac{1}{2^k}$ converges).  So, the Comparison Test says that
    %   $\displaystyle \sum \frac{5 + 3\sin k}{2^k}$ converges, too. 
    %   \end{solution}
    %     \probonly{\vfill}
        
        
    % \end{enumerate}     




  \end{enumerate}

\end{document}